% !PRINTED 588
\weberpage{588}

Since $A$ and $D$ are without common divisors, we have, for a variable $x$ and a constant $h$:
\[
h\,x^{m-1} A\!\left(\frac{1}{x}\right) = D(x),
\]
and, on replacing $x$ by $1/x$:
\[
h\,A(x) = x^{m-1} D\!\left(\frac{1}{x}\right),
\]
and by (16):
\[
h^2 = (-1)^{cx}.
\]
Thus
\weq{17}{D(x) = \varepsilon\,x^{m-1} A\!\left(\frac{1}{x}\right),}
where $\varepsilon = \pm 1$ or $= \pm i$. By (16) this depends on whether $cx$ is even or odd\footnote{From § 138 it follows that $i$ is contained in the field $\mathfrak{L}$; (17) is therefore not in contradiction with Theorem 2.}.

If we therefore use the notation (10) again, but now under the hypothesis that $A$ and $D$ are without common divisors, we obtain:
\weq{18}{\varepsilon\,\boldsymbol{s}(\mu u) = \frac{A_1\boldsymbol{s} + A_3\boldsymbol{s}^3 + \cdots A_{m-2}\boldsymbol{s}^{m-2} + \boldsymbol{s}^m}{1 + A_{m-2}\boldsymbol{s}^2 + \cdots A_3\boldsymbol{s}^{m-3} + A_1\boldsymbol{s}^{m-1}},}
and the $A_1, A_3, \ldots, A_{m-1}$ are \gesperrt{integers of the field} $\mathfrak{L}$, for they are the symmetric fundamental functions of the integers (13).
% [source sic: the prose reaches A_{m-1}, whereas equation (18) ends the unknown sequence at A_{m-2}]

In particular, if one sets $u=0$, it follows that:
\weq{19}{A_1 = \varepsilon\mu,}
and from this:
\weq{20}{\varepsilon\mu = \prod\nolimits^{\varrho} \boldsymbol{s}\!\left(\frac{2\varrho}{\mu}\right).}
The $\boldsymbol{s}\!\left(\dfrac{2\varrho}{\mu}\right)$ are therefore divisors of $\mu$.

From (10) it follows that:
\weq{21}{\boldsymbol{s}A(\boldsymbol{s}) - \boldsymbol{s}(\mu u)D(\boldsymbol{s}) = 0,}
and this, when $\boldsymbol{s}(\mu u)$ is regarded as given, is an equation for $\boldsymbol{s}(u)$ of degree $m$, whose roots are:
\[
\boldsymbol{s}(u),\ \boldsymbol{s}\!\left(u+\frac{2\varrho}{\mu}\right).
\]
%CONT>>
